\section{An exact coercive identity}\label{cert:section}

We retain the fixed positive time, positive constant step, and notation of
Lemma~\ref{shape:ibp}.  We will prove
\[
 \E_f[z^2-12vw^2+6v^4]\geq\frac1{1000}\E_fv^4.
\]
The expression need not be positive pointwise.  Subtracting a spatial
derivative with zero expectation preserves its expectation; the calculation
below decomposes the resulting expression into nonnegative terms.  The variables introduced at the end of the
previous section satisfy
\begin{equation}\label{cert:signs}
 R,K,v,B,j,L,S\geq0,\qquad b=B+2v\geq0,
 \qquad Rb^2=w^2.
\end{equation}
In particular, the signs of $B$ and $S$ record the backward inequalities
$H\leq0$ and $H_x\geq0$.

For a polynomial $F$ in $R,K,v,b,j$, set
\begin{equation}\label{cert:operators}
 \mathcal D F=r\partial_xF,\qquad
 \mathcal W(F)=\mathcal DF+(v+R-K)F.
\end{equation}
The derivative rules in~\eqref{shape:D-rules} determine $\mathcal W$
algebraically.  Its analytic meaning is
\begin{equation}\label{cert:divergence}
 f\mathcal W(F)=\partial_x(rfF),\qquad
 \E_f\mathcal W(F)=0.
\end{equation}
Lemma~\ref{shape:ibp} supplies the integrability and vanishing boundary
terms for every polynomial used here.  Thus it suffices to find $F$ for
which subtracting $\mathcal W(F)$ from the curvature expression leaves a
positive multiple of $v^4$ and other nonnegative terms.

\subsection{Choosing the derivative correction}

We first arrange the terms containing $z$.  Set
\begin{equation}\label{cert:constants-first}
 (A,C,k,T_0,g_0,h)
 =\left(\frac1{13},\frac{53}{24},\frac{19}{16},\frac{73}{48},
         \frac{2331}{10000},\frac{87}{2500}\right)
\end{equation}
and define
\begin{equation}\label{cert:QT}
 Q=Avb+Cv^2+kKv-2kRb+T_0Rv,
 \qquad T=g_0Rv+hjv.
\end{equation}
The terms $(z+Q)^2$ and $ST$ will be nonnegative: the first is a square,
and $S,T\geq0$.  Since $S=z+bv-2Rb$, their terms linear in $z$ have
total coefficient $2Q+T$.

The only derivative rule containing $z$ is
$\mathcal Db=-z-bv$.  Consequently the coefficient of $z$ in
$\mathcal W(F)$ is $-\partial_bF$.  We therefore choose $F=F_2$ so that
$\partial_bF_2=2Q+T$.  Set
\begin{equation}\label{cert:constants-second}
 (x_0,y,m,n_0)
 =\left(\frac{4663}{5000},\frac{6853}{10000},
          -\frac{1051}{400},-\frac{1073}{500}\right).
\end{equation}
The following polynomial has the required derivative:
\begin{equation}\label{cert:polynomials}
\begin{aligned}
 F_2={}&(Av-2kR)b^2\\
       &+b\bigl(2Cv^2+2kKv+(g_0+2T_0)Rv+hjv\bigr)\\
       &+x_0v^2j+yv^3+mKv^2+n_0Rv^2.
\end{aligned}
\end{equation}
Its first two lines are the required antiderivative in $b$; the last line
is independent of $b$ and will be used in the remaining terms.  Directly,
\[
 \partial_bF_2
 =2Avb-4kRb+2Cv^2+2kKv+(g_0+2T_0)Rv+hjv
 =2Q+T.
\]
It follows that
\[
 z^2-12Rvb^2+6v^4-(z+Q)^2-ST-\mathcal W(F_2)
\]
has no $z$ dependence.  This is the reason for the choice of the derivative
correction.

The remaining calculation uses one further favorable term from the forward
score inequality.  Let $\mathcal W_0$ be the polynomial rule obtained from
$\mathcal W$ by omitting $-L$ in $\mathcal Dj=-jv+Kb-L$.  Since
\begin{equation}\label{cert:L-correction}
 \partial_jF_2=x_0v^2+hbv,
 \qquad
 \mathcal W(F_2)=\mathcal W_0(F_2)-L(x_0v^2+hbv),
\end{equation}
replacing $\mathcal W$ by $\mathcal W_0$ leaves the nonnegative term
$L(x_0v^2+hbv)$ to be retained in the final identity.  We now have only a
polynomial in $R,K,v,b,j$ to decompose.

\subsection{The nonnegative remainder}

Define the five polynomials
\begin{equation}\label{cert:squares}
\begin{aligned}
 P_0&=Rj(b-7v)^2,\\
 P_1&=RB\left(b-\frac73v\right)^2,\\
 P_2&=Kv\left(b-\frac92v\right)^2,\\
 P_3&=vB\left(K+\frac14R-\frac17v\right)^2,\\
 P_4&=\left(vb-\frac{14}{5}v^2-\frac{13}{4}Kv+\frac43Rv\right)^2,
\end{aligned}
\end{equation}
with positive coefficients
\begin{equation}\label{cert:constants}
 (p_0,p_1,p_2,p_3,p_4)
 =\left(\frac{337}{10000},\frac{183}{2500},\frac{411}{10000},
         \frac{11461}{5000},\frac{647}{10000}\right).
\end{equation}
Let $P_{\mathrm{rem}}$ be the polynomial in $R,K,v,B,j$ whose nonzero
coefficients are listed in Table~\ref{cert:table}.  All its variables and
coefficients are nonnegative.

\begin{lemma}[Rational polynomial identity]\label{cert:identity-lemma}
With these definitions, $B=b-2v$, and $S=z+bv-2Rb$, one has the polynomial
identity over $\mathbb Q$
\begin{equation}\label{cert:identity}
\begin{aligned}
 z^2-12Rvb^2+6v^4
  ={}&(z+Q)^2+ST+\sum_{i=0}^4p_iP_i+\mathcal W(F_2)\\
     &+L(x_0v^2+hbv)+P_{\mathrm{rem}}.
\end{aligned}
\end{equation}
\end{lemma}

\begin{proof}
Equation~\eqref{cert:L-correction} reduces the assertion to the same identity
with $\mathcal W_0$ and without the displayed $L$ term.  The construction of
$F_2$ above cancels every term containing $z$.  Expand the residual
\[
 z^2-12Rvb^2+6v^4-(z+Q)^2-ST
       -\mathcal W_0(F_2)-\sum_{i=0}^4p_iP_i
\]
using~\eqref{shape:D-rules}, substitute $b=B+2v$, and collect monomials in
$R,K,v,B,j$.  The coefficients are exactly those of
Table~\ref{cert:table}, and every unlisted coefficient is zero.  This proves
the identity by rational arithmetic.
\end{proof}

\begin{table}[htbp]
\centering
\renewcommand{\arraystretch}{1.16}
\begin{tabular}{c|r@{\qquad}c|r}
Monomial & Coefficient & Monomial & Coefficient\\ \hline
$R^2v^2$ & $223/160000$ & $R^2vB$ & $4007/240000$\\
$R^2B^2$ & $1097/120000$ & $RKv^2$ & $83/80000$\\
$RKvB$ & $3343/24000$ & $Rv^3$ & $16259/23400000$\\
$Rv^2B$ & $30179/5460000$ & $Rv^2j$ & $26/625$\\
$RvB^2$ & $107/15000$ & $RvBj$ & $3/2500$\\
$RB^3$ & $121/32500$ & $RB^2j$ & $11/10000$\\
$K^2v^2$ & $579/20000$ & $K^2vB$ & $207/2500$\\
$Kv^3$ & $911/975000$ & $Kv^2B$ & $349/420000$\\
$Kv^2j$ & $5011/5000$ & $KvB^2$ & $133/130000$\\
$KvBj$ & $87/2500$ & $v^4$ & $1661507/1521000000$\\
$v^3B$ & $633799/621075000$ & $v^2B^2$ & $10657/1690000$
\end{tabular}
\caption{All 22 nonzero coefficients of $P_{\mathrm{rem}}$; every unlisted
coefficient is zero.  Each listed coefficient is positive; the coefficient
of $v^4$ gives the coercive lower bound.}
\label{cert:table}
\end{table}

\begin{proposition}[Integrated coercivity]\label{cert:coercivity}
Suppose $r,v,w,z,s$ and the variables $R,K,b,j,L,B,S$ are related as in
Lemma~\ref{shape:ibp}, with the derivative rules~\eqref{shape:D-rules}
and signs~\eqref{cert:signs}.  Assume their quotients extend across zero.
Let $f$ be a probability density satisfying $f_x/f=r-s$ such that every
term of~\eqref{cert:identity} is integrable and $rfF_2$ tends to zero at both
spatial infinities.  Then
\begin{equation}\label{cert:inequality}
 \E_f[z^2-12vw^2+6v^4]
 \geq\frac{1661507}{1521000000}\E_fv^4
 \geq\frac1{1000}\E_fv^4.
\end{equation}
\end{proposition}

\begin{proof}
Integrating~\eqref{cert:identity} removes $\mathcal W(F_2)$ by
\eqref{cert:divergence}.  Every remaining term on the right is nonnegative.
Indeed, the prefactors of the five squares $P_i$ are $Rj,RB,Kv,vB,1$;
$T=v(g_0R+hj)\geq0$; and $L(x_0v^2+hbv)\geq0$.
The table gives the stronger bound
\[
 P_{\mathrm{rem}}\geq\frac{1661507}{1521000000}v^4,
 \qquad
 \frac{1661507}{1521000000}-\frac1{1000}
       =\frac{140507}{1521000000}>0.
\]
Using $Rb^2=w^2$ proves~\eqref{cert:inequality}.
\end{proof}

For the smooth finite-step data of Section~\ref{shape:section},
Lemma~\ref{shape:ibp} verifies all the hypotheses of this proposition.
Although the proof uses the density score and its ratios, the resulting
inequality involves only $v,w,z$.  This is the form that passes to an
integrable order parameter.

